% Job posting: https://picovoice.ai/careers/deep-learning-researcher/
% =====================================================================
%  Cover Letter - Nishal Stanislaus Silva, PhD
%  Target: Picovoice - Deep Learning Researcher (Vancouver)
%  Variant: V3 (Edge / Embedded), research-forward
% =====================================================================

\documentclass[11pt]{article}

\usepackage[left=0.8in,top=0.7in,right=0.8in,bottom=0.7in]{geometry}
\usepackage{enumitem}
\usepackage[hidelinks]{hyperref}
\usepackage{setspace}
\usepackage{parskip}
\usepackage[en-GB]{datetime2}
\usepackage[T1]{fontenc}
\usepackage{newtxtext,newtxmath}
\ifdefined\pdfgentounicode
  \input{glyphtounicode}
  \pdfgentounicode=1
\fi

\setlength{\parindent}{0pt}
\setstretch{1.05}
\pagenumbering{gobble}

\newcommand{\clName}{Nishal Stanislaus Silva}
\newcommand{\clEmail}{nishal.silva@hotmail.com}
\newcommand{\clPhone}{+1 613 200 2030}
\newcommand{\clCity}{Toronto, ON}
\newcommand{\clSite}{https://nishal.xyz}
\newcommand{\clLinkedIn}{https://www.linkedin.com/in/nishal-silva}
\newcommand{\clStatus}{Canadian Permanent Resident, Eligible to work in Canada without sponsorship}

\newcommand{\clTagline}{ML Research Engineer | Real-Time \& Edge Inference | Audio, Sensor \& Time-Series}
\newcommand{\clRole}{Deep Learning Researcher}
\newcommand{\clCompany}{Picovoice}
\newcommand{\clTeam}{the research team}
\newcommand{\clAddressee}{Dear Hiring Manager,}

\begin{document}
\pagestyle{empty}

\begin{center}
    {\LARGE \scshape \textbf{\clName}}{\large \textit{, PhD}}
    \\[1pt]
    {\small \clTagline}
    \\[1pt]
    {\footnotesize\textit{\clStatus}}
    \\[1pt]
    {\small
    \href{mailto:\clEmail}{\clEmail}
    \,\,|\,\, \clPhone
    \,\,|\,\, \clCity
    \,\,|\,\, \href{\clSite}{nishal.xyz}
    \,\,|\,\, \href{\clLinkedIn}{linkedin.com/in/nishal-silva}}
\end{center}

\rule{\linewidth}{0.4pt}\vspace{0.6em}

\today

\textbf{Re: \clRole{} -- \clCompany{}}

\clAddressee

I am applying for the Deep Learning Researcher role at Picovoice. I am a deep learning researcher and engineer with a PhD from the University of Trento whose thesis is the same problem Picovoice solves: real-time audio understanding that runs entirely on the device, under tight latency and compute budgets. My published work runs audio pattern detection in 14 ms on a Raspberry Pi 4 at 31.4 percent CPU, F1 0.76, which is 74 times faster than the dynamic time warping baseline it replaced (JAES 2026).

The role asks for someone who understands deep learning theory and its limits, can build complex and extensible software in Python and C, and can solve vague problems on a tight schedule. My PhD and postdoc were exactly that: I designed the model architectures, feature representations and training regimes, built the baseline and ablation suites that showed which choices actually mattered against latency and CPU, and shipped the result as a running system on embedded hardware. I wrote and maintain Nebula, an open-source C++17 audio-feature library with per-feature ARM latency benchmarks, and I work daily with PyTorch and statistical evaluation. My GPU work has been training rather than CUDA kernel development, so that is where I would ramp, and I would ramp fast given the C and profiling background behind it.

What draws me to Picovoice is the conviction that voice AI should run offline, on the user's hardware, with no cloud in the loop. I have spent years making models small and fast enough to earn that place, and I would bring a benchmark and dataset discipline, four open datasets and 70+ contributors' worth, to how new models are measured before they ship. In the first months I would contribute to model efficiency and evaluation on the on-device speech and audio work.

I am a Canadian permanent resident based in Toronto, eligible to work without sponsorship, available immediately, and glad to relocate to Vancouver for this role. I would welcome the opportunity to discuss how my experience in real-time and on-device ML can support \clTeam{} at Picovoice. My portfolio, publications and open datasets are at \href{\clSite}{nishal.xyz}.

\vspace{12pt}
Sincerely,\\[2pt]
\textbf{\clName}, PhD

\end{document}
